\section{Practice and Current Discipline}

This section is a study aid describing how the received sources say the
Rosary is well prayed, and what the current indulgence discipline
provides. It issues no devotional obligations: under MC 55 the faithful
are ``serenely free'' in the Rosary's regard, and nothing here
substitutes for the liturgy, for pastoral guidance, or for a
confessor's counsel.

\subsection{Praying it well: the papal suggestions}

Drawn from RVM 28--37 and MC 46--50, all as suggestions carrying their
authors' pastoral authority and no more:

\begin{itemize}
\item \textbf{Announce the mystery}, perhaps with an icon or image; the
announcement opens ``a scenario on which to focus our attention'' (RVM
29).
\item \textbf{Let Scripture speak}: a related biblical passage, long or
short, may follow the announcement---``No other words can ever match
the efficacy of the inspired word'' (RVM 30). The Rosary presupposes
and promotes \term{lectio divina}; it does not replace it (RVM 29).
\item \textbf{Keep silence} briefly after the word, before the vocal
prayer (RVM 31).
\item \textbf{Weigh the words}: the Our Father grounds each decade in
the Father; the ten Aves carry the contemplation, the name of Jesus
their hinge---in public recitation some regions add a clause naming the
mystery; the doxology is a summit, not a formality (RVM 32--34; MC
46).
\item \textbf{Pace and tone}: ``a quiet rhythm and a lingering pace''
(MC 47); recitation ``grave and suppliant'' in the Lord's Prayer,
``lyrical and full of praise'' in the Aves, ``full of adoration'' in
the doxology (MC 50).
\item \textbf{Conclude in the Church}: prayer for the Pope's
intentions, and a Marian antiphon or litany as the crowning of the
inner journey (RVM 37); a concluding collect follows local custom (RVM
35).
\item \textbf{Family recitation} is the setting the modern popes urge
most: the family Rosary gathers parents and children before the image
of the Virgin, unites them with the absent and the dead (Pius XII),
and makes the family that prays it ``reproduce something of the
atmosphere of the household of Nazareth'' (RVM 41); the Rosary is
``one of the best and most efficacious prayers in common'' of the
family (MC 54).\endnote{\work{Ingruentium malorum} 12--14; RVM 41--42;
MC 52--54.}
\end{itemize}

Two cautions bound the practice. The Rosary is not to be recited during
the celebration of the liturgy (MC 48), and its fruits are invitations
to conversion, not guarantees: no recitation count purchases an
outcome, and the beads are not a charm (RVM 28). Reported private
revelations---Lourdes and Fatima above all---have commended this
prayer, and the Church has recognized their influence; but private
revelation ``is not the same as public revelation,'' binds no one to
believe, and did not originate the Rosary, which was centuries old
before either apparition.\endnote{RVM 7 with its note 11 (``private
revelations are not the same as public revelation, which is binding on
the whole Church''); the Fatima events postdate by three and a half
centuries the form stabilized under Pius V. Apparition judgments and
reported messages have their own competent processes and records and
are outside this reference's scope.}

\subsection{The current indulgence grant}

\begin{disciplinebox}{Enchiridion Indulgentiarum, 4th ed.\ (1999),
concession 17 \S 1, as of 2026-07-25}
A \textbf{plenary indulgence} is granted to the member of the faithful
who devoutly recites the Marian Rosary (1) in a church or oratory, or
in a family, a religious community, or an association of the faithful,
and in general when several gather for an honest purpose; or (2) who
devoutly unites himself to the recitation of the prayer by the Supreme
Pontiff broadcast by television or radio. In other circumstances the
indulgence is \textbf{partial}. For the plenary grant the Enchiridion
further provides: (a) the recitation of a third part of the Rosary
(five decades) suffices, but the five decades must be recited
continuously; (b) devout meditation on the mysteries must be added to
the vocal prayer; and (c) in public recitation the mysteries must be
announced according to approved local custom, while in private
recitation it suffices to join meditation on the mysteries to the
vocal prayer.
\end{disciplinebox}

The general norms cap and condition every such grant: a plenary
indulgence can be acquired only once a day; it requires, besides
exclusion of all attachment to even venial sin, the performance of the
indulgenced work and the three conditions of sacramental confession,
Eucharistic communion, and prayer for the Supreme Pontiff's
intentions; and one who for a reasonable cause completes only part of a
divisible work (the Rosary is the Enchiridion's own example, divisible
into decades) acquires a partial indulgence for the part
performed.\endnote{\work{Enchiridion Indulgentiarum}, 4th ed., norm 18
\S 1, norm 20 \S 1, and praenotanda 6--7, checked in the Latin web
presentation cited in section~1 on 2026-07-25. The prose statement of
concession 17 \S 1 above is an editorial summary of the Latin, not an
official translation.} Indulgence discipline is mutable law of the
Apostolic Penitentiary: the statement above is true of the fourth
edition as checked on 2026-07-25, historical grants (such as Leo
XIII's 1883 concessions) are superseded, and a reader or future editor
must verify the current Enchiridion---and, where a vernacular prayer
text's approval matters for a grant, the competent authority's approval
of that translation---before relying on any of it.\endnote{See the
terminal appendix for the as-of boundary. The prayer texts printed in
section~1 are a received historical witness printed for study; this
reference makes no claim that that particular printing is an approved
text for any indulgence-related purpose.}

\subsection{What the Rosary is not}

The boundaries, gathered: the Rosary is not liturgy and not
sacramental anamnesis, though it lives from the liturgy and leads back
to it (MC 48; RVM 13); it is not an obligation, though it is warmly and
continuously commended (MC 55); it is not a relic of pre-conciliar
piety, the Council's liturgical primacy notwithstanding (RVM 4); it is
not un-Christological---its Marian mode carries a wholly
Christ-centered content (RVM 1, 18); it is not a magical practice (RVM
28); and it adds nothing to public Revelation, which is complete in
Christ---the mysteries propose to memory and love what the Gospel
already contains (CCC 65--67; RVM 29).\endnote{CCC 65--67 on the
completeness of Revelation in Christ; the remaining loci as cited. On
Mary's subordinate, participated cooperation---the doctrinal frame of
every Marian devotion---see LG 60--62 and CCC 970: her maternal
intercession ``in no way obscures or diminishes the unique mediation
of Christ, but rather shows its power'' (LG 60 in the rendering quoted
at RVM 15; the Vatican conciliar English reads ``in no wise obscures
or diminishes this unique mediation of Christ, but rather shows His
power'').}
